\chapter{Derivators}
\label{chap:derivators}

In this chapter, we collect all our previous results and interpret them to say that a cofibration category has a canonically associated {\it left Heller derivator}, and dually a fibration category has a canonically associated {\it right Heller derivator}. 

Derivators are best thought of as an axiomatization of homotopy (co)limits. They were introduced by Alexandre Grothendieck in his manuscripts \cite{Grothendieck1}, \cite{Grothendieck2}. Don Anderson \cite{Anderson2} and Alex Heller \cite{Heller} present alternative definitions of derivators. For an elementary introduction to the theory of derivators, the reader is invited to refer to Georges Maltsiniotis \cite{Maltsiniotis1}.

Grothendieck's original definition had derivators as contravariant on both 1- and 2- cells; we will refer to this type of derivators as {\it Grothendieck derivators}. Heller's definition had derivators as contravariant on 1- cells but covariant on 2-cells; we will refer to them as {\it Heller derivators}, although the axioms we use are set up a little bit differently than in \cite{Heller}. The axioms are arranged so that Grothendieck and Heller derivators correspond to each other by duality.

We will recall the basic definitions we need. Fix a pair of universes ${\cat U} \subset {\cat U}^{'}$. We denote $2{\cat Cat}$ the 2-category of ${\cat U}$-small categories, and $2CAT$ the 2-category of ${\cat U}^{'}$-small categories. We denote $\initialcat$ the initial category (having an empty set of objects), and $\terminalcat$ the terminal category (having one object and the identity of that object as the only map). We write $p_{\cat D} \colon {\cat D} \ra \terminalcat$ for the functor to the terminal category and $e_{d} \colon \terminalcat \ra {\cat D}$ for the functor that embeds the terminal category $\terminalcat$ as the object $d \eps {\cat D}$. 

For a 2-category ${\cat C}$, we denote ${\cat C}^{1-op}$ (${\cat C}^{2-op}$, resp. ${\cat C}^{1,2-op}$) the 2-category obtained reversing the direction of the 1-cells (resp. 2-cells, resp. both 1- and 2-cells) of ${\cat C}$.

We will denote by ${\cat Dia}$ any 2-full 2-subcategory of $2{\cat Cat}$, with the property that its objects (viewed as small categories) include the set of finite direct categories, and are stable under the following category operations:
\begin{enumerate}
\item
small disjoint sums of categories
\item
finite products of categories
\item
stable under taking overcategories and undercategories, i.e. if $f \colon {\cat D}_1$ $\ra$ ${\cat D}_2$ is a functor with ${\cat D}_1, {\cat D}_2 \eps {\cat Dia}$ then $(f \downarrow d_2)$, $(d_2 \downarrow f)$ $\eps {\cat Dia}$ for any object $d_2 \eps {\cat D_2}$
\end{enumerate}
In particular from (1) and (2), we will assume that $\initialcat$, $\terminalcat$ $\eps {\cat Dia}$.

For example, we can take ${\cat Dia}$ to be
\begin{enumerate}
\item
the entire $2{\cat Cat}$, or
\item
the 2-full subcategory $2{\cat FinDirCat}$ whose objects are small disjoint sums of finite direct categories
\end{enumerate}

\section{Prederivators}
\label{sec:preder}

\begin{defn}
\label{defn:hellerpreder}
A {\it Heller prederivator}\index{prederivator} of domain ${\cat Dia}$ is a pseudo 2-functor ${\mathbb D} \colon {\cat Dia}^{1-op} \ra  2CAT$. 
\end{defn}

For any morphism $u \colon {\cat D}_1 \ra {\cat D}_2$ in ${\cat Dia}$ we denote $u^* = {\mathbb D}u \colon {\mathbb D}({\cat D}_2) \ra {\mathbb D}({\cat D}_1)$, and for any natural map $\alpha : u \Ra v$ in ${\cat Dia}$ we denote $\alpha^* = {\mathbb D}\alpha : u^* \ra v^*$.

A Heller prederivator is {\it strict} if it is strict as a 2-functor.

For example, any category with weak equivalences $({\cat M}, {\cat W})$ gives rise to a strict Heller prederivator ${\mathbb D}_{({\cat M}, {\cat W})}$ of domain $2{\cat Cat}$ defined by ${\mathbb D}_{({\cat M}, {\cat W})} ({\cat D}) = {\bf ho}{\cat M}^{{\cat D}}$, where the homotopy category is taken with respect to pointwise weak equivalences ${\cat W}^{\cat D}$. Any natural map $\alpha \colon u \Ra v$ yields a natural map $\alpha^* \colon u^* \Ra v^*$
\begin{center}
  $\xymatrix{
    {\cat D}_1 \ar@/^1pc/[r]^u _\,="a" \ar@/_1pc/[r]_v ^\,="b" & {\cat D_2} && {\bf ho}{\cat M}^{{\cat D}_1}  && {\bf ho}{\cat M}^{{\cat D}_2} \ar@/^1pc/[ll]^{v^*} _\,="c" \ar@/_1pc/[ll]_{u^*} ^\,="d"
    \ar@{=>} "a";"b" ^{\alpha}
    \ar@{=>} "d";"c" ^{\alpha^*}
    }$
\end{center}
where $\alpha^*$ is defined on components as $X_{\alpha d_1}: X_{vd_1} \ra X_{ud_1}$, for diagrams $X \eps {\cat M}^{{\cat D}_1}$.

\section{Derivators}
\label{sec:der}
Suppose that $u \colon {\cat D}_1 \ra {\cat D}_2$ is a functor in ${\cat Dia}$.  A left adjoint for $u^* \colon {\mathbb D}({\cat D}_2) \ra {\mathbb D}({\cat D}_1)$, if it exists, is denoted $u_{!} \colon {\mathbb D}({\cat D}_1) \ra {\mathbb D}({\cat D}_2)$. A right adjoint for $u^*$, if it exists, is denoted $u_{*} \colon {\mathbb D}({\cat D}_1) \ra {\mathbb D}({\cat D}_2)$. If the functors $u_!$ or $u_*$ exist, they are only defined up to unique isomorphism.

Suppose we have a diagram in ${\cat Dia}$ with $\phi \colon vf \Ra gu$ a natural map
\begin{equation}
\label{eqn:2catsquareder}
\xymatrix {
    {\cat D}_1 \ar[r]^u \ar[d]_f & 
    {\cat D}_2 \ar[d]^g _\,="b" \\
    {\cat D}_3 \ar[r]_v ^\,="a" & 
    {\cat D}_4 \ultwocell<\omit>{<0>\phi}
  }
\end{equation}
We apply the 2-functor ${\mathbb D}$ and we get a natural map $\phi^* \colon f^*v^* \Ra u^*g^*$. If the left adjoints $u_!$ and $v_!$ exist, the natural map $\phi^*$ yields by adjunction the {\it cobase change} natural map denoted
\begin{equation}
\label{eqn:2catsquaredermap}
\phi_! \colon u_! f^* \Ra g^* v_!
\end{equation}
If $\phi_!$ is an isomorphism for a choice of left adjoints $u_!, v_!$, then $\phi_!$ is an isomorphism for any such choice of right adjoints $u_!, v_!$. Dually, if we assume that the right adjoints $f^*$ and $g^*$ exist, the {\it base change} morphism associated to $\phi$ is
\begin{equation}
\label{eqn:2catsquaredermapop}
\phi_* \colon v^* g_* \Ra f_* u^*
\end{equation}
and if $\phi_*$ is an isomorphism for a choice of $f_*, g_*$ then $\phi_*$ is an isomorphism for any such choice.

Given a functor $u \colon {\cat D}_1 \ra {\cat D_2}$ and an object $d_2 \eps {\cat D_2}$, in the standard over 2-category diagram of $u$ at $d_2$ (see (\ref{eqn:stdover2cat}))
\begin{equation}
\label{eqn:stdover2catbasechangedgm}
\xymatrix {
    (u \downarrow d_2) \ar[r]^-{p_{(u \downarrow d_2)}} \ar[d]_{i_{u, d_2}} & 
    \terminalcat \ar[d]^{e_{d_2}} _\,="b" \\
    {\cat D}_1 \ar[r]_u ^\,="a" & 
    {\cat D}_2 \ultwocell<\omit>{<0>_{{\phi}_{u, d_2} \;\;\;\;\;\;}}
  }
\end{equation}

In \eqref{eqn:stdover2catbasechangedgm}, if the left adjoints $u_!$ and $p_{(u \downarrow d_2)!}$ exist, the associated cobase change morphism is denoted
\begin{equation}
\label{eqn:stdover2catbasechange}
({\phi}_{u, d_2})_! \colon p_{(u \downarrow d_2)!}i_{u, d_2}^* \Ra e_{d_2}^*u_!
\end{equation}

Dually, for an object $d_2 \eps {\cat D_2}$, in the standard under 2-category diagram of $u$ at $d_2$ (see (\ref{eqn:stdunder2cat}))
\begin{equation}
\label{eqn:stdover2catcobasechangedgm}
\xymatrix {
    (d_2 \downarrow u) \ar[r]^-{i_{d_2, u}} \ar[d]_{p_{(d_2 \downarrow u)}}& 
    {\cat D}_1 \ar[d]^u _\,="b" \\
    \terminalcat \ar[r]_{e_{d_2}} ^\,="a" & 
    {\cat D}_2 \ultwocell<\omit>{<0>_{{\phi}_{d_2, u} \;\;\;\;\;\;}}
  }
\end{equation}

In \eqref{eqn:stdover2catcobasechangedgm}, if the right adjoints $u^*$ and $p_{(d_2 \downarrow u)}^*$ exist, the associated base change morphism is denoted
\begin{equation}
\label{eqn:stdunder2catcobasechange}
({\phi}_{d_2, u})^* \colon e_{d_2}^*u_* \Ra p_{(d_2 \downarrow u)*}i_{d_2, u}^*
\end{equation}

\index{derivator}
\index{derivator!Heller derivator}
\begin{defn}[Heller derivators]
\label{defn:hellerderivator}
\mbox{} \\
A Heller prederivator ${\mathbb D} \colon {\cat Dia}^{1-op} \ra  2CAT$ is a {\it left Heller derivator} if it satisfies the following axioms:
\begin{description}
\item[Der1]
For any set of small categories ${\cat D}_k$, $k \eps K$, the functor
\begin{center}
$\xymatrix {
    {\mathbb D}( {\Sum}_{k \eps K} {\cat D}_k) 
    \ar[rr]^-{(i_k^*)_{k \eps K}} &&
    {\times}_{k \eps K} {\mathbb D}({\cat D}_k)
  }$
\end{center}
is an equivalence of categories, where $i_{k} \colon {\cat D}_{k} \ra {\Sum}_{k^{'} \eps K} {\cat D}_{k^{'}}$ denotes the component inclusions for $k \eps K$.

\item[Der2]
For any ${\cat D}$ in ${\cat Dia}$, the family of functors $e_{d}^* \colon {\mathbb D}({\cat D}) \ra {\mathbb D}(\terminalcat)$ for all objects $d \eps {\cat D}$ is conservative.

\item[Der3l]
For any $u \colon {\cat D}_1 \ra {\cat D}_2$ in ${\cat Dia}$, the functor $u^* \colon {\mathbb D}({\cat D}_2) \ra {\mathbb D}({\cat D}_1)$ admits a left adjoint $u_! \colon {\mathbb D}({\cat D}_1) \ra {\mathbb D}({\cat D}_2)$

\item[Der4l]
For any $u \colon {\cat D}_1 \ra {\cat D}_2$ in ${\cat Dia}$ and any object $d_2 \eps {\cat D}_2$, the cobase change morphism $({\phi}_{u, d_2})_! \colon p_{(u \downarrow d_2)!}i_{u, d_2}^* \Ra e_{d_2}^*u_!$ of (\ref{eqn:stdover2catbasechange}) associated to the standard over 2-category diagram (\ref{eqn:stdover2cat}) is an isomorphism.
\end{description}

${\mathbb D}$ is a {\it right Heller derivator} if it satisfies axioms Der1 and Der2 above and:
\begin{description}
\item[Der3r]
For any $u \colon {\cat D}_1 \ra {\cat D}_2$ in ${\cat Dia}$, the functor $u^* \colon {\mathbb D}({\cat D}_2) \ra {\mathbb D}({\cat D}_1)$ admits a right adjoint $u_* \colon {\mathbb D}({\cat D}_1) \ra {\mathbb D}({\cat D}_2)$

\item[Der4r]
For any $u \colon {\cat D}_1 \ra {\cat D}_2$ in ${\cat Dia}$ and any object $d_2 \eps {\cat D}_2$, the base change morphism $({\phi}_{d_2, u})^* \colon e_{d_2}^*u_* \Ra p_{(d_2 \downarrow u)+}i_{d_2, u}^*$ of (\ref{eqn:stdunder2catcobasechange}) associated to the standard under 2-category diagram (\ref{eqn:stdunder2cat}) is an isomorphism.
\end{description}

${\mathbb D}$ is a {\it Heller derivator} if it is both a left and a right Heller derivator. 
\end{defn}

\index{derivator}
\index{derivator!Grothendieck derivator}
\begin{defn}[Grothendieck derivators]
\label{defn:grothderivator}
A 2-functor ${\mathbb D} \colon {\cat Dia}^{1,2-op} \ra  2CAT$ is a left (right, two-sided) {\it Grothendieck derivator} if the functor ${\cat Dia}^{1-op} \ra  2CAT$, ${\cat D} \mapsto {\mathbb D}({\cat D}^{op})$ is a right (left, two-sided) Heller derivator.
\end{defn}

Notice that for a left Grothendieck derivator the map $u^* = {\mathbb D}(u)$ by definition admits a {\it right} adjoint $u_*$, and for a right Grothendieck derivator ${\mathbb D}(u)$ admits a {\it left} adjoint $u_!$. This observation should be kept in mind when reading \cite{Grothendieck1}, \cite{Grothendieck2}, \cite{Maltsiniotis1}, \cite{Cisinski2}, \cite{Cisinski3}, \cite{Cisinski4}. 

We have allowed for pseudo-functors ${\mathbb D}$ rather than just strict 2-functors in the definition of Heller, resp. Grotherndieck derivators. Grothendieck and Heller had originally just used strict 2-functors ${\mathbb D}$ in the derivator definition, but their arguments apply just as well for pseudo-functors.

The axioms for a Grothendieck derivator that we propose are stronger than the ones originally introduced by Grothendieck in that we require Der1 to hold for arbitrary small sums of categories rather than just finite sums of categories. The axioms for a Heller derivator that we propose are weaker than the axioms of a Heller {\it homotopy theory}, as defined in \cite{Heller}, in that our Heller derivators don't satisfy Heller's axiom H2 (same as our Der5 below).

We will use the term {\it derivator} to refer to Heller derivators, when no confusion is possible.

For ${\cat D}_1$ and ${\cat D}_2$ in ${\cat Dia}$, we define the functor
\begin{center}
$\mathrm{dgm}_{_{{\cat D}_1, {\cat D}_2}} \colon \mathbb{D}({\cat D}_1 \times {\cat D}_2) \ra \mathbb{D}({\cat D}_2)^{{\cat D}_1}$ \\
$(\mathrm{dgm}_{_{{\cat D}_1, {\cat D}_2}} X)_{d_1} = i_{d_1}^* X$ \\
$(\mathrm{dgm}_{_{{\cat D}_1, {\cat D}_2}} X)_{f} = i_{f}^* X$
\end{center}
for any object $d_1$ and map $f \colon d_1 \ra d^{'}_1$ of ${\cat D}_1$, where $i_{d_1} \colon {\cat D}_2 \ra {\cat D}_1 \times {\cat D}_2$ is the functor $d_2 \mapsto (d_1, d_2)$ and $i_f \colon i_{d_1} \Ra i_{d^{'}_1}$ is the natural map given on components by $(f, 1_{{\cat D}_1})$.

\index{derivator!strong}
\begin{defn}[Strong Heller derivators]
\label{defn:hellerstrongderivator}
\mbox{} \\
A left (resp. right, two-sided) Heller derivator ${\mathbb D} \colon {\cat Dia}^{1-op} \ra  2CAT$ is {\it strong} if it satisfies the additional axiom Der5
\begin{description}
\item[Der5]
For any finite, free category ${\cat D}_1$ and category ${\cat D}_2$ in ${\cat Dia}$, the functor $\mathrm{dgm}_{_{{\cat D}_1, {\cat D}_2}}$ is full and essentially surjective.
\end{description}
\end{defn}

Note that in the literature the axiom Der5 is sometimes stated in the weaker form
\begin{description}
\item[Der5w]
For any category ${\cat D}_2$ in ${\cat Dia}$ the functor $\mathrm{dgm}_{_{I, {\cat D}_2}}$ is full and essentially surjective, where $I$ denotes the category $0 \ra 1$, with two objects and one non-identity map.
\end{description}

\comment{
\index{derivator!pointed}
\begin{defn}[Pointed Heller derivators]
\label{defn:hellerpointedderivator}
\mbox{} \\
A left (resp. right, two-sided) Heller derivator ${\mathbb D} \colon {\cat Dia}^{1-op} \ra  2CAT$ is {\it pointed} if it satisfies the additional axiom Der6l (resp. Der6r, resp. Der6l and Der6r).
\begin{description}
\item[Der6l]
For any closed embedding $u \colon {\cat D}_1 \ra {\cat D}_2$ in ${\cat Dia}$, the functor $u_!$ admits a left adjoint denoted $u^?$
\item[Der6r]
For any open embedding $u \colon {\cat D}_1 \ra {\cat D}_2$ in ${\cat Dia}$, the functor $u_*$ admits a right adjoint denoted $u^!$
\end{description}
\end{defn}
} %end comment

\begin{thm}
\label{thm:hellerderishtpytheory}
A Heller prederivator of domain $2{\cat Cat}$ is a {\it homotopy theory} in the sense of \cite{Heller} iff it is a strong, strict Heller derivator in the sense of \definitionref{defn:hellerderivator}.
\end{thm}

\begin{proof}
This is a consequence of the results of \cite{Heller}. 
\end{proof}

In fact, a careful reading of \cite{Heller} shows that most proofs stated there for one-sided homotopy theories are true for one-sided {\it strong} Heller derivators. However we suspect that the theory developed in \cite{Heller} can be reworked from the axioms of a one-sided Heller derivator alone. 

\section{Derivability of cofibration categories}
\label{sec:derivability}
Let us turn back to the Heller prederivator ${\mathbb D}_{({\cat M}, {\cat W})}$ associated to a category with weak equivalences $({\cat M}, {\cat W})$ and introduce the following

\index{derivable category}
\begin{defn}
\label{defn:derivablecat}
A category with weak equivalences $({\cat M}, {\cat W})$ is {\it left Heller derivable} (resp. strongly left Heller derivable, right Heller derivable, etc) over ${\cat Dia}$ if the prederivator ${\mathbb D}_{({\cat M}, {\cat W})}$ is a left (left strong, right etc.) Heller derivator over ${\cat Dia}$.
\end{defn}

A left Heller derivable (resp. strongly left Heller derivable, etc.) category by convention is left Heller derivable over the entire $2{\cat Cat}$. 

With the definitions of derivators and of derivable categories at hand, we can state

\begin{thm}[Cisinski]
\label{thm:cofcatderivator}
\mbox{}
\begin{enumerate}
\item
Any ABC cofibration category is strongly left Heller derivable.
\item
Any ABC fibration category is strongly right Heller derivable.
\item
Any ABC model category is strongly Heller derivable.
\end{enumerate}
\end{thm}

\begin{proof}
Part (2) is dual to (1) and part (3) is a consequence of (1) and (2). Therefore it suffices to prove part (1). 

Part (1) is proved by verifying the derivator axioms as follows:
\begin{enumerate}
\item[-]
Der1 is proved by \lemmaref{lem:cofcatsumsofdiagrams}
\item[-]
Der2 by \theoremref{thm:homotcatconservative}
\item[-]
Der3r by \theoremref{thm:genholim}
\item[-]
Der4r by \theoremref{thm:computerelhocolim}
\item[-]
Der5 in \theoremref{thm:generalpointwisecofstructure} and \theoremref{thm:freedirectdiagramrealization}
\end{enumerate}
\end{proof}

As a corollary of \theoremref{thm:cofcatderivator} and \propositionref{prop:relclosedmodcat} we obtain

\begin{thm}
\label{thm:modelcatderivator}
Any Quillen model category is strongly Heller derivable. $\square$
\end{thm}

In particular, all the results of \cite{Heller} apply to ABC model categories (and therefore to Quillen model categories).

We should mention that our terminology for derivable categories is different than the one used by Denis-Charles Cisinski. His paper \cite{Cisinski2} calls left derivable categories what we call F1-F4 Anderson-Brown-Cisinski fibration categories. The naming anomaly is explained by the fact that \cite{Cisinski2} is written in the language of Grothendieck derivators, and a left Grothendieck derivator corresponds to a right Heller derivator.

\comment{
To illustrate the power of \theoremref{thm:cofcatderivator}, we end with the following application.

\begin{thm}
\label{thm:homotcofinalhomotocolim}
Suppose that $u \colon {\cat D}_1 \ra {\cat D}_2$ is a functor of small categories. Then:
\begin{enumerate}
\item
$u$ is homotopy right cofinal iff for any cofibration category ${\cat M}$ and any diagram $X \eps {\cat M}^{{\cat D}_2}$ the natural map 
\begin{center}
${\bf L}\colim^{{\cat D}_1} \; {\bf ho}u^*X \ra {\bf L}\colim^{{\cat D}_2}X$
\end{center}
is an isomorphism in ${\bf ho}{\cat M}$.
\item
$u$ is homotopy left cofinal iff for any fibration category ${\cat M}$ and any diagram $X \eps {\cat M}^{{\cat D}_2}$ the natural map 
\begin{center}
${\bf R}\lim^{{\cat D}_2}X \ra {\bf R}\lim^{{\cat D}_1} \; {\bf ho}u^*X$ 
\end{center}
is an isomorphism in ${\bf ho}{\cat M}$.
\end{enumerate}
\end{thm}
} %comment

\section{Odds and ends}
\label{sec:oddsandends}
The author has not been able to complete his research project as thoroughly as he would have liked, but it may be worthwhile to at least state how this body of work may be improved and extended. First of all, it is proved by Denis-Charles Cisinski that 

\begin{thm}[\cite{Cisinski2}]
\label{thm:cofcatderivator2}
\mbox{}
\begin{enumerate}
\item
Any CF1-CF4 cofibration category is strongly left Heller derivable over $2{\cat FinDirCat}$. 
\item
Any F1-F4 fibration category is strongly right Heller derivable over $2{\cat FinDirCat}$.
\end{enumerate}
\end{thm}

Indeed, we have modeled our proofs in a very large part on the arguments of \cite{Cisinski2}.

Second, consider a weaker replacement of axiom CF6.

\begin{description}
\item[CF6']
For any sequence of weak equivalence maps $A_0 \la B_0 \ra A_1 \la B_1 \ra A_2 ...$ such that each $B_n$ is cofibrant, and such that $B_0 \ra A_0$ and all $B_n \Sum B_{n+1} \ra A_n$ are cofibrations, we have that 
\[
A_0 \ra \colim (A_0 \la B_0 \ra A_1 \la B_1 \ra A_2 ...)
\]
is a weak equivalence.
\end{description}

The colimit in CF6' can be shown to always exist using a modified version of \lemmaref{lem:colimpushout}. If  $A^{'}_0 \la B^{'}_0 \ra A^{'}_1 ...$ is another sequence of weak equivalence maps such that each $B_n$ is cofibrant, and such that $B^{'}_0 \ra A^{'}_0$ and all $B^{'}_n \Sum B^{'}_{n+1} \ra A^{'}_n$ are cofibrations, and if $a_n \colon A_n \ra A^{'}_n$, $b_n \colon B_n \ra B^{'}_n$ form a {\it diagram} map, then using a modified version of \theoremref{thm:colimdirect} one can show that the induced map
\[
\colim (A_0 \la B_0 \ra A_1 \la B_1 \ra A_2 ...) \ra \colim (A^{'}_0 \la B^{'}_0 \ra A^{'}_1 \la B^{'}_1 \ra A^{'}_2 ...)
\]
is a weak equivalence.

We conjecture that a precofibration category $({\cat M}, {\cat W}, {\cat Cof})$ satisfying CF5 and CF6' is strongly left Heller derivable. It is an open question actually if axiom CF6' is needed at all - so an even stronger conjecture would be that axioms CF1-CF5 imply strong left Heller derivability.

The dual axiom for fibration categories is 

\begin{description}
\item[F6']
For any sequence of weak equivalence maps $A_0 \ra B_0 \la A_1 \ra B_1 \la A_2 ...$ such that each $B_n$ is fibrant, and such that $A_0 \ra B_0$ and all $A_{n+1} \ra B_n \times B_{n+1}$ are fibrations, we have that 
\[
\lim (A_0 \ra B_0 \la A_1 \ra B_1 \la A_2 ...) \ra A_0
\]
is a weak equivalence.
\end{description}

The dual weak conjecture is that a prefibration category $({\cat M}, {\cat W}, {\cat Fib})$ satisfying F5 and F6' is strongly right Heller derivable. The dual strong conjecture is that a prefibration category satisfying F5 is strongly right Heller derivable.
